\chapter{Tools for assessing the scientific potentials of detector configurations \label{ch:data_analysis}}

To compute the various data analysis quantities quoted in this report, we use the classic gravitational-wave data analysis computations based on matched filtering with a stationary Gaussian Likelihood function. Our main goal in gravitational-wave data analysis is to get the posterior distribution $\pi(\vec{\theta}|D)=p(D|\vec{\theta})p(\vec{\theta})/p(D)$, where $p(D|\vec{\theta})$ is the Likelihood ($\mathcal{L}(\vec{\theta})$) of matching the data given the model parameters, $p(\vec{\theta})$ is the prior probability, and $p(D)$ is the evidence which is an integral of the numerator over all of parameter space. 

Given time-domain dimensionless strain data d(t) and theoretical template h(t), the log of the Likelihood function is given by,
\begin{equation}\label{eq:log_like}
    \ln{\mathcal{L}\sim-\frac{1}{2}\langle d - h | d-h\rangle=-\frac{1}{2}\left(\langle d|d\rangle + \langle h|h\rangle - 2 \langle d|h\rangle\right)},
\end{equation}
where $\langle a | b \rangle$ is the noise-weighted inner product of the two time-domain signals $a(t)$ and $b(t)$. This is computed in the frequency domain:
\begin{equation}\label{eq:inner_product}
    \langle a|b \rangle = 4\Re \int_{-\infty}^\infty \frac{\tilde{a}(f)^*\tilde{b}(f)}{S_n(f)}df,
\end{equation}
where $\tilde{a}(f)$ is the Fourier transform of $a(t)$ and $S_n(f)$ is the total noise budget during observation. The noise budget typically includes two components: the instrument noise and any confusion noise from stochastic signals~\cite{Robson:2018ifk,Karnesis:2021tsh}; $S_n(f)$ represents the direct sum of all these components at frequency $f$. When discretized into a sum, we add the numerical noise normalization component necessary for fitting stochastic signals: $-\sum_i\log{(2\pi S_n(f_i))}$~\cite{Katz:2024oqg}.

Most of the computations we look at involve computing the signal-to-noise ratio (SNR) of a given signal. The SNR is given by $\sqrt{\langle h |h\rangle}$. When we go beyond the optimal SNR of a possible source to determine uncertainty in its parameters, we compute this in two ways: i) estimates from a Fisher information matrix formalism and ii) numerical generation of posterior distributions using Markov Chain Monte Carlo (MCMC) and similar methods~\cite{Karnesis:2023ras,Katz:2024oqg}. The information matrix $M(\vec{\theta})$ for parameters $\vec{\theta}$ with template $h(t, \vec{\theta})$ is given by the inner product of partial derivatives of $h(t, \vec{\theta})$ with respect to each parameter:
\begin{equation}
    M_{ij}(\vec{\theta}) = \langle h_i(\vec{\theta})|h_j(\vec{\theta})\rangle,
\end{equation}
where $h_i(\vec{\theta})$ is the partial derivative of $h(t, \vec{\theta})$ with respect to the parameter $\{\vec{\theta}\}_i$. Inverting $M$ gives the theoretical covariance matrix at $\vec{\theta}$. This Fisher information matrix is theoretical in nature and relies on a high-SNR assumption on the source of interest~\cite{Vallisneri:2007ev}. The information matrix can be correct for many cases when trying to understand theoretical uncertainties. When something more accurate is needed, we use Markov Chain Monte Carlo (MCMC) and similar methods. The technique of MCMC involves numerically exploring the Likelihood surface especially in the vicinity of the true parameters (so-called ``typical set''). MCMC can provide more accurate estimates for the posterior distributions and, therefore, the parameter uncertainties with the Fisher information matrix representing the theoretical limit for the parameter covariances. MCMC can be very computationally expensive, so it is mainly used where other tools are unable to properly compute accurate results.

When examining LISA, we have to make a key distinction from that of current ground-based observations: the complexity of the LISA response function. The LISA response to a gravitational-wave signal is both time- and frequency-dependent. It contains two parts: i) projection of the waveforms onto the links of the detector and ii) combination of those link projections with Time Delay Interferometry (TDI)~\cite{Tinto:2020fcc} into TDI observables. The LISA response function is naturally described in the time-domain where due to the time-dependence of the delays between spacecraft. However, we can more easily describe TDI mathematically in the frequency domain, where we can represent the effect of the response as a time- and frequency-dependent transfer function based on the time-frequency correspondence, $t(f)$~\cite{Marsat:2018oam}:
\begin{equation}
    \tilde{h}^\text{X,Y,Z}(f)=\mathcal{T}^\text{X,Y,Z}(t(f), f)\tilde{h}^\text{SSB}(f),
\end{equation}
where $\mathcal{T}^\text{X,Y,Z}$ is the transfer function, $h^\text{SSB}$ is the intrinsic waveform located at the solar system barycenter, and  $h^\text{X,Y,Z}$ could also be written ($\tilde{X}(f), \tilde{Y}(f), \tilde{Z}(f)$) representing the individual X, Y, and Z TDI channels. X, Y, and Z are correlated Michelson channels. Using a linear transformation, we can change to channels A, E, T that are uncorrelated under assumptions of equal and constant LISA armlengths~\cite{Prince:2002hp}, which is sufficient for SNR and data analysis investigations for future detector predictions. The stochastic confusion and noise components are also projected onto this TDI basis. Therefore, Eq.~\ref{eq:inner_product} becomes a sum over three independent channels A,E, and T: $\langle a|b\rangle=\sum_\text{i=X,Y,Z}\langle a^i|b^i\rangle$. 

To compute the SNR for a specific source-type, you apply the specific waveform to each source type as $h^\text{X,Y,Z}$. Generating time-domain waveforms can be very slow for these predictive computations. Therefore, source-by-source we build fast computations. The source waveforms themselves can all be found in the gravitational-wave community. However, adjustments for fast frequency-domain LISA response functions for different sources and detectors can be nuanced. The ratio of frequency derivative due to the response ($\dot{f}_\text{LISA}$ versus the frequency derivative of the source ($\dot{f}_\text{GW}$) can affect how frequency-domain response functions are computed. At ($\dot{f}_\text{LISA}<<\dot{f}_\text{GW}$), the fast frequency domain response can be computed as a function of frequency and taking the spacecraft positions to be a function of $t(f)$ where $t(f)$ is monotonic and increasing~\cite{Marsat:2020rtl,Katz:2022yqe}. At the opposite end ($\dot{f}_\text{LISA}>>\dot{f}_\text{GW}$ and $\dot{f}_\text{GW}T_\text{obs} << \text{mHz}$), you can have multiple points in time represented at each frequency as LISA constellation movement doppler shifts the signal back and forth. In this case, you can compute the response at very large timesteps and heterodyne around the carrier frequency of the gravitational wave~\cite{Cornish:2007if}.
